% 가상 신청서에 configs/costs_klm.yaml을 적용한 예. 수집 자료의 평가값이 아니다.
% src/formbench/editing/cost.py align(): FORMBENCH_COST_UNITS=klm, lang=ko, 허용오차 12pt.
% 목표: 성명=홍길동, 생년월일=1990. 1. 1., 주소=서울시 중구 (본문에서 설명).
% 원본 23.96초: 값 세 칸 입력. 행 누락 31.35초: 행 추가 + 주소 항목 + 값 세 칸 입력.
% 이미지 52.09초: 이미지 삭제 1.58 + 표 추가 50.51(표 삽입 8.65 + 여섯 칸 입력).
% 서론의 직관 예시(교수 피드백 2026-10-19): 외관의 유사성과 편집 시간이 어긋남을 보인다. 막대 높이는 반올림 전 편집 시간에 비례, 표시는 초.
% 서론 둘째 문단 바로 뒤에 오도록 [H]로 두고(사용자 요청 2026-10-20), 1페이지 왼쪽 단의 남은 높이에 맞게 낮게 그린다:
% 막대는 값 칸 아래에 두고, "편집 시간"은 표의 항목 칸 아래에 행 머리로 둔다.
\begin{figure}[H]
\centering
\begin{tikzpicture}[
  lc/.style={draw,minimum width=12mm,minimum height=3.8mm,inner sep=0pt,font=\scriptsize},
  vc/.style={draw,minimum width=15mm,minimum height=3.8mm,inner sep=0pt,font=\scriptsize},
  head/.style={font=\footnotesize,anchor=south,align=center,inner sep=1pt},
  value/.style={font=\footnotesize\bfseries,anchor=south,inner sep=1.5pt}]
\newcommand{\blankform}[1]{%
  \node[lc] at (#1+0.60,-0.19) {성명}; \node[vc] at (#1+1.95,-0.19) {};
  \node[lc] at (#1+0.60,-0.57) {생년월일}; \node[vc] at (#1+1.95,-0.57) {};}
\node[head] at (1.35,0.04) {원본\\{\scriptsize 기입만 필요}};
\node[head] at (4.55,0.04) {복원본 A\\{\scriptsize 행 하나 누락}};
\node[head] at (7.75,0.04) {복원본 B\\{\scriptsize 표를 이미지로 삽입}};
\blankform{0}
\node[lc] at (0.60,-0.95) {주소}; \node[vc] at (1.95,-0.95) {};
\blankform{3.2}
\draw[dashed,black!55] (3.2,-0.76) rectangle (5.9,-1.14) (4.4,-0.76) -- (4.4,-1.14);
\begin{scope}[every node/.append style={fill=black!15,draw=black!55,text=black!55}]
  \blankform{6.4}
  \node[lc] at (7.00,-0.95) {주소}; \node[vc] at (8.35,-0.95) {};
\end{scope}
% 공통의 0 기준선 위에서 길이만으로 시간 차이를 읽는다. 막대는 각 표의 값 칸 아래.
\draw[black!50] (1.30,-2.28) -- (9.10,-2.28);
\node[font=\footnotesize,anchor=south,inner sep=1pt] at (0.60,-2.28) {편집 시간};
\foreach \x/\ratio/\display in {1.95/1/24초,5.15/1.30843/31초,8.35/2.17404/52초} {
  \fill[black!60] (\x-0.40,-2.28) rectangle (\x+0.40,-2.28+0.32*\ratio);
  \node[value] at (\x,-2.28+0.32*\ratio) {\display};
}
\end{tikzpicture}
\caption{외관의 유사성과 편집 시간이 어긋나는 예}
\label{fig:cost-example}
\end{figure}
